\appendix

\chapter{Code Listings}
\label{app:code}

\section{Channel Generation (3GPP UMa)}

\begin{lstlisting}[language=Python, caption={3GPP TR 38.901 UMa Channel Model Implementation}]
import numpy as np

def generate_uma_channel(N, cell_radius=5000, fc=3.5e9):
    """
    Generate 3GPP UMa channel for N users.
    
    Parameters:
    - N: Number of users
    - cell_radius: Cell radius in meters
    - fc: Carrier frequency in Hz
    
    Returns:
    - h: Channel gains (N,)
    - positions: User positions (N, 2)
    """
    # Random user positions (uniform in circle)
    r = cell_radius * np.sqrt(np.random.rand(N))
    theta = 2 * np.pi * np.random.rand(N)
    x = r * np.cos(theta)
    y = r * np.sin(theta)
    positions = np.stack([x, y], axis=1)
    
    # 2D distance to BS at origin
    d2D = np.sqrt(x**2 + y**2)
    
    # UE height: 1.5 m, BS height: 25 m
    hUT = 1.5
    hBS = 25
    
    # 3D distance
    d3D = np.sqrt(d2D**2 + (hBS - hUT)**2)
    
    # Clamp to valid range [10, 5000] meters
    d3D = np.clip(d3D, 10, 5000)
    
    # LOS probability (Table 7.4.2-1)
    PL_LOS = (18 * np.log10(d3D) + 28.0 + 
              20 * np.log10(fc / 1e9))
    
    # Calculate LOS probability
    p_LOS = np.where(
        d2D <= 18,
        1.0,
        18 / d2D + np.exp(-d2D / 63) * (1 - 18 / d2D)
    )
    
    # Random LOS/NLOS assignment
    is_LOS = np.random.rand(N) < p_LOS
    
    # NLOS path loss (Table 7.4.1-1)
    PL_NLOS = (161.04 - 7.1 * np.log10(1.5) + 
               7.5 * np.log10(40) - 
               (24.37 - 3.7 * (hBS / 25)**2) * 
               np.log10(hBS) + 
               (43.42 - 3.1 * np.log10(hBS)) * 
               (np.log10(d3D) - 3) + 
               20 * np.log10(fc / 1e9) - 
               (3.2 * (np.log10(11.75 * 1.5))**2 - 4.97))
    
    # Use minimum of LOS and NLOS for each user
    PL = np.where(is_LOS, PL_LOS, 
                  np.maximum(PL_LOS, PL_NLOS))
    
    # Shadow fading (log-normal)
    sigma_SF = np.where(is_LOS, 4, 6)  # dB
    SF = np.random.randn(N) * sigma_SF
    
    # Total path loss in dB
    PL_total = PL + SF
    
    # Rayleigh fading (exponential distribution)
    rayleigh = np.random.exponential(1.0, N)
    rayleigh_dB = 10 * np.log10(rayleigh)
    
    # Channel gain in dB
    h_dB = -PL_total + rayleigh_dB
    
    # Convert to linear scale
    h = 10 ** (h_dB / 10)
    
    return h, positions, d2D, theta
\end{lstlisting}

\section{GNN Model Architecture}

\begin{lstlisting}[language=Python, caption={Multi-Task GraphSAGE Model}]
import torch
import torch.nn as nn
import torch.nn.functional as F
from torch_geometric.nn import SAGEConv

class PairPowerGNN(nn.Module):
    def __init__(self, in_channels=5, hidden_channels=128, 
                 num_layers=3, dropout=0.3):
        super().__init__()
        
        # GraphSAGE encoder
        self.convs = nn.ModuleList()
        self.convs.append(SAGEConv(in_channels, hidden_channels))
        for _ in range(num_layers - 1):
            self.convs.append(
                SAGEConv(hidden_channels, hidden_channels)
            )
        
        self.dropout = dropout
        
        # Decoder heads
        # 1. Pairing classifier
        self.pair_mlp = nn.Sequential(
            nn.Linear(2 * hidden_channels, hidden_channels),
            nn.ReLU(),
            nn.Dropout(dropout),
            nn.Linear(hidden_channels, 1)
        )
        
        # 2. Sum-rate regressor
        self.sumrate_mlp = nn.Sequential(
            nn.Linear(2 * hidden_channels, hidden_channels),
            nn.ReLU(),
            nn.Dropout(dropout),
            nn.Linear(hidden_channels, 1)
        )
        
        # 3. Power allocator
        self.power_mlp = nn.Sequential(
            nn.Linear(2 * hidden_channels, hidden_channels),
            nn.ReLU(),
            nn.Dropout(dropout),
            nn.Linear(hidden_channels, 2)  # [P_i, P_j]
        )
    
    def forward(self, x, edge_index):
        # Encoder: message passing
        for i, conv in enumerate(self.convs):
            x = conv(x, edge_index)
            if i < len(self.convs) - 1:
                x = F.relu(x)
                x = F.dropout(x, p=self.dropout, 
                             training=self.training)
        
        # Node embeddings: x has shape [N, hidden_channels]
        return x
    
    def decode(self, z, edge_index):
        # z: node embeddings [N, hidden_channels]
        # edge_index: [2, E] edges to predict
        
        src, dst = edge_index
        edge_emb = torch.cat([z[src], z[dst]], dim=-1)
        
        # Three predictions per edge
        pair_logits = self.pair_mlp(edge_emb).squeeze(-1)
        sumrate = self.sumrate_mlp(edge_emb).squeeze(-1)
        power = self.power_mlp(edge_emb)  # [E, 2]
        
        return pair_logits, sumrate, power
\end{lstlisting}

\section{Training Loop}

\begin{lstlisting}[language=Python, caption={Multi-Task Training Procedure}]
def train_epoch(model, loader, optimizer, device):
    model.train()
    total_loss = 0
    
    for batch in loader:
        batch = batch.to(device)
        optimizer.zero_grad()
        
        # Forward pass: encode nodes
        z = model(batch.x, batch.edge_index)
        
        # Decode on all edges
        pair_logits, sumrate_pred, power_pred = \
            model.decode(z, batch.edge_label_index)
        
        # Multi-task loss
        # 1. Pairing BCE
        pair_loss = F.binary_cross_entropy_with_logits(
            pair_logits, 
            batch.edge_label.float()
        )
        
        # 2. Sum-rate MAE (only on positive pairs)
        pos_mask = batch.edge_label == 1
        if pos_mask.sum() > 0:
            sumrate_loss = F.l1_loss(
                sumrate_pred[pos_mask],
                batch.sumrate_label[pos_mask]
            )
        else:
            sumrate_loss = torch.tensor(0.0, device=device)
        
        # 3. Power MAE (only on positive pairs)
        if pos_mask.sum() > 0:
            power_loss = F.l1_loss(
                power_pred[pos_mask],
                batch.power_label[pos_mask]
            )
        else:
            power_loss = torch.tensor(0.0, device=device)
        
        # Weighted combination
        loss = (1.0 * pair_loss + 
                0.5 * sumrate_loss + 
                0.3 * power_loss)
        
        loss.backward()
        torch.nn.utils.clip_grad_norm_(
            model.parameters(), 1.0
        )
        optimizer.step()
        
        total_loss += loss.item()
    
    return total_loss / len(loader)
\end{lstlisting}

\section{Inference and Matching}

\begin{lstlisting}[language=Python, caption={Greedy Pairing with SIC Verification}]
def infer_pairing(model, data, device, sic_threshold=8.0):
    model.eval()
    with torch.no_grad():
        data = data.to(device)
        
        # Encode nodes
        z = model(data.x, data.edge_index)
        
        # Get all pairwise scores
        N = data.num_nodes
        edge_index = torch.combinations(
            torch.arange(N), r=2
        ).T.to(device)
        
        pair_logits, sumrate_pred, _ = \
            model.decode(z, edge_index)
        
        # Convert logits to probabilities
        pair_probs = torch.sigmoid(pair_logits)
    
    # Greedy matching on CPU
    scores = pair_probs.cpu().numpy()
    edges = edge_index.T.cpu().numpy()
    h = data.x[:, 0].cpu().numpy()  # channel gains
    
    # Sort edges by score descending
    sorted_idx = np.argsort(-scores)
    
    paired = set()
    final_pairs = []
    
    for idx in sorted_idx:
        i, j = edges[idx]
        
        # Skip if already paired
        if i in paired or j in paired:
            continue
        
        # SIC constraint: |h_i - h_j| >= threshold dB
        h_diff_dB = 10 * np.log10(h[i] / h[j])
        if abs(h_diff_dB) >= sic_threshold:
            final_pairs.append((int(i), int(j)))
            paired.add(i)
            paired.add(j)
    
    return final_pairs
\end{lstlisting}

\section{Dataset Generation}

\begin{lstlisting}[language=Python, caption={Training Dataset Creation}]
import pandas as pd
from tqdm import tqdm

def generate_dataset(num_scenarios=50000, save_path='dataset.csv'):
    data = []
    
    for scenario_id in tqdm(range(num_scenarios)):
        # Random user count
        N = np.random.choice([4, 6, 8, 10, 12])
        
        # Generate UMa channel
        h, pos, d, theta = generate_uma_channel(N)
        
        # Compute optimal pairing via Blossom
        pairs = blossom_matching(h, pos)
        
        # Store scenario
        for user_i in range(N):
            data.append({
                'scenario_id': scenario_id,
                'user_id': user_i,
                'h': h[user_i],
                'x': pos[user_i, 0],
                'y': pos[user_i, 1],
                'd': d[user_i],
                'theta': theta[user_i],
                'N': N
            })
        
        # Store pairing labels
        for i, j in pairs:
            # Symmetric edges
            data.append({
                'scenario_id': scenario_id,
                'edge': (i, j),
                'label': 1,
                'sumrate': compute_sumrate(h[i], h[j]),
                'power': optimize_power(h[i], h[j])
            })
    
    df = pd.DataFrame(data)
    df.to_csv(save_path, index=False)
    print(f"Saved {num_scenarios} scenarios to {save_path}")

# Run generation
generate_dataset(50000, 'data/raw/dataset.csv')
\end{lstlisting}

\section{Evaluation Script}

\begin{lstlisting}[language=Python, caption={Complete Evaluation on Test Set}]
def evaluate_test_set(model, test_loader, device):
    model.eval()
    
    all_pair_probs = []
    all_pair_labels = []
    all_sumrate_preds = []
    all_sumrate_labels = []
    
    total_throughput_gnn = 0
    total_throughput_blossom = 0
    
    with torch.no_grad():
        for batch in test_loader:
            batch = batch.to(device)
            
            # Encode
            z = model(batch.x, batch.edge_index)
            
            # Decode
            pair_logits, sumrate_pred, _ = \
                model.decode(z, batch.edge_label_index)
            
            pair_probs = torch.sigmoid(pair_logits)
            
            # Collect metrics
            all_pair_probs.append(pair_probs.cpu())
            all_pair_labels.append(batch.edge_label.cpu())
            
            pos_mask = batch.edge_label == 1
            if pos_mask.sum() > 0:
                all_sumrate_preds.append(
                    sumrate_pred[pos_mask].cpu()
                )
                all_sumrate_labels.append(
                    batch.sumrate_label[pos_mask].cpu()
                )
            
            # Throughput comparison
            gnn_pairs = infer_pairing(model, batch, device)
            gnn_throughput = sum(
                compute_sumrate(batch.x[i, 0], batch.x[j, 0])
                for i, j in gnn_pairs
            )
            
            blossom_throughput = batch.throughput_opt.item()
            
            total_throughput_gnn += gnn_throughput
            total_throughput_blossom += blossom_throughput
    
    # Classification metrics
    pair_probs = torch.cat(all_pair_probs)
    pair_labels = torch.cat(all_pair_labels)
    
    from sklearn.metrics import roc_auc_score
    auc = roc_auc_score(pair_labels, pair_probs)
    
    # Regression metrics
    sumrate_preds = torch.cat(all_sumrate_preds)
    sumrate_labels = torch.cat(all_sumrate_labels)
    mae = F.l1_loss(sumrate_preds, sumrate_labels).item()
    
    # Throughput ratio
    throughput_ratio = (total_throughput_gnn / 
                       total_throughput_blossom * 100)
    
    print(f"Test AUC: {auc:.4f}")
    print(f"Sum-Rate MAE: {mae:.4f} bps/Hz")
    print(f"Throughput vs Blossom: {throughput_ratio:.2f}%")
    
    return {
        'auc': auc,
        'mae': mae,
        'throughput_ratio': throughput_ratio
    }
\end{lstlisting}
